\documentclass[11pt,letterpaper]{article}
\usepackage[T1]{fontenc}
\usepackage{lmodern,amsmath,amssymb,amsthm,mathtools}
\usepackage[margin=0.88in]{geometry}
\usepackage{booktabs,longtable,array,microtype,xcolor}
\usepackage{enumitem,fancyhdr,hyperref,xurl}
\hypersetup{colorlinks=true,linkcolor=blue!45!black,urlcolor=blue!45!black,citecolor=blue!45!black,
 pdftitle={A single norm for factorial and digit certificates},
 pdfauthor={AI-assisted research prepared for Vladimir Reshetnikov}}
\newtheorem{theorem}{Theorem}[section]
\newtheorem{lemma}[theorem]{Lemma}
\newtheorem{proposition}[theorem]{Proposition}
\newtheorem{corollary}[theorem]{Corollary}
\theoremstyle{remark}\newtheorem{remark}[theorem]{Remark}
\newcommand{\Z}{\mathbb Z}\newcommand{\N}{\mathbb N}
\newcommand{\Con}{\operatorname{Con}}
\newcommand{\G}{\mathcal G}
\newcommand{\Pnew}{P_{245}}
\setlength{\emergencystretch}{2em}
\setlist{itemsep=3pt,topsep=5pt}
\pagestyle{fancy}\fancyhf{}
\fancyhead[L]{\small Diophantine consistency}
\fancyhead[R]{\small Joint factorial and digit norm}
\fancyfoot[C]{\thepage}\setlength{\headheight}{14pt}
\title{A single norm for factorial and digit certificates\\[5pt]
\large A direct quartic for ZFC consistency in 245 integer variables}
\author{AI-assisted research prepared for Vladimir Reshetnikov\\
\small ProveIt research program}
\date{25 October 2026}
\begin{document}
\maketitle
\begin{abstract}
We give a fixed polynomial of total degree four in 245 unrestricted integer
variables whose integer unsolvability is provably equivalent in ZFC to
\(\Con(\mathrm{ZFC})\). Its complete specification is the sum of squares of
181 explicitly listed quadratic residuals. Its collected expansion has
5,304 nonzero monomials and coefficient height 128,446. Against the latest
inspected direct construction in 252 integers, this saves seven variables,
one residual, and 36 monomials while preserving degree and coefficient height.
The mathematical change makes one common norm certify both the low digits
of binomial expansions and a nearest-integer factorial identity. It removes
all three initial sign certificates. An odd radix multiplier and a careful
order of interpretation recover the signs afterwards. A separate argument
handles the zero factorial candidate, where the usual coefficient bound
fails. The construction uses the repository's numerical proof calculus
C222; it does not encode a K215 computation or instantiate Jones's universal
polynomial. The new proof and finite algebraic certificates compose with
identified inherited arithmetic and refutation calibrations. The result is
an improvement of this direct family, not a global simplicity theorem or a
new end-to-end proof-assistant certificate. A separate companion equation
has ten strictly positive variables, exact degree 454, and 78 binary
arithmetic operations; its proof separates two square conditions only after
recovering the main Pell index.
\end{abstract}
\clearpage
\tableofcontents
\clearpage

\section{The exact result and the comparison boundary}

Let \(X=(X_0,\ldots,X_{244})\) range over \(\Z^{245}\). Define
\begin{equation}\label{eq:main}
 \boxed{\quad \Pnew(X)=\sum_{i=0}^{180}R_i(X)^2=0.\quad}
\end{equation}
Every coefficient and every index in the 181 residuals \(R_i\) is printed in
Appendix~\ref{app:residuals}. The same literal list is supplied in
\path{direct/equation/equation.json}; the independently recomputed collected
expansion is in \path{direct/equation/expanded_polynomial.json}. These files
define an ordinary polynomial over \(\Z\). All abbreviations in the proof
below describe the construction; no variable exponent, factorial, binomial
operator, hidden parameter, or unspecified universal polynomial occurs in
the delivered polynomial.

\begin{theorem}\label{thm:main}
For the fixed polynomial \(\Pnew\) specified above, the arithmetic and
syntactic-calibration arguments in this article and the retained source
proofs establish
\[
 \mathrm{ZFC}\ \vdash\
 \left(\forall X\in\Z^{245}\;\Pnew(X)\ne0\right)
 \ \longleftrightarrow\ \Con(\mathrm{ZFC}).
\]
Equivalently, ZFC proves that an integer zero exists if and only if there is
a finite ZFC refutation.
\end{theorem}

Here consistency means syntactic consistency for the fixed effective ZFC
proof presentation used by the inherited calibration. The theorem establishes
an equivalence, and assumes neither consistency nor inconsistency in its
construction. All coefficients are generated by finite arithmetic independent
of the existence of a refutation.

The comparison snapshot is the exact public ProveIt commit
\par\noindent\texttt{3ecd6a6cf53d47519228328567f45fb7ee9781d9}.
All eight archives then
present in \path{docs/incoming} were inspected. The immediate baseline is the
\emph{Common Remainders and Radix 47} archive, specifically its direct
252-variable construction \cite{repo,common}. The normalized source, its
proof closure, source identities, licenses, and calibration interface files
are retained under \path{direct/inherited/normalized252}.

\begin{center}
\begin{tabular}{lrrr}
\toprule
Measure & Inspected parent & This construction & Change\\
\midrule
Unrestricted integer variables & 252 & 245 & $-7$\\
Quadratic residuals & 182 & 181 & $-1$\\
Total degree & 4 & 4 & 0\\
Collected monomials & 5,340 & 5,304 & $-36$\\
Coefficient height & 128,446 & 128,446 & 0\\
\bottomrule
\end{tabular}
\end{center}

The coefficient height is the largest absolute coefficient after collection,
not the size of a factored expression. The largest residual coefficient is
256. The expansion has constant coefficient 2,348. The comparison concerns
literal equations, with integer domains kept fixed.

\paragraph{What simplicity means here.}
The repository also contains equations with far fewer variables and much
larger degree or fixed coefficient payloads. In particular, its compact
nine-positive and ten-integer constructions are not displaced by a
245-variable quartic \cite{forms,sun}. Even among quartics, the native
coefficient-encoding family has different tradeoffs. The present contribution
improves the direct C222 family: small literal coefficients, degree four,
and no enormous packed coefficient tuple. There is no known global
minimum being certified here. The new idea is the removal of the initial
sign layer through a common arithmetic certificate.

\section{The inherited interface and the new equations}

\subsection{What is retained}

The direct source represents a finite numerical C222 derivation by CRT
coordinates, a factorial-based family of coprime moduli, and a collection
of power requests. A signed Pell construction verifies these requests.
The source has eighteen power requests and fourteen independent principal
Pell pairs, with an additional pair obtained by one recurrence step.
Its square frame and squared semantic anchor bound and orient the classified
Pell indices before front arithmetic signs are used. The retained proof
\path{ORIENTATION_FIRST_PROOF.md} is the relevant sign-order argument;
\path{STRICT_BUDGET_PROOF.md} fixes the amplification exponent.

For the new proof we use the following interface, supplied by those retained
ordinary proofs.

\begin{enumerate}
\item The horizon is \(t=16+h_0^2\), so \(t\ge16\).
The square frame bounds the absolute values of the actual bases, requested
indices, and outputs by a number \(W\ge258\). Its signed Pell center
\(A\) satisfies \(|A|>W^W\).
\item Every principal pair is the canonical Pell polynomial pair at its
requested signed integer index. This conclusion comes from bounded residue
errors and a separate squared anchor, not from the signs of \(F,B,q\).
\item For a requested base \(n\) with \(|n|\ge2\), the canonical power
congruence forces its index \(k\) to be nonnegative and its output to equal
\(n^k\). For base 1 and a known nonnegative index it forces output 1.
We justify the signed-base extension in Section~\ref{sec:signed}.
\item The local C222 arithmetic, CRT decoding, amplification bound, and
finite proof-flow argument apply once the genuine factorial, powers, and
binomial digit intervals have been recovered.
\end{enumerate}

These are inherited mathematical inputs with retained proof texts, not new
claims of a machine-checked compilation of the whole equation. In particular,
the signed index bound is available without any of the three initial sign
guards removed below. This dependency is essential.

\subsection{Notation at the arithmetic interface}

The letters in this subsection denote semantic registers of the source,
before its exact eliminations and translations. They need not be coordinates
with the same names in the final residual list. Write
\[
 V=(F^2+1)^t,\qquad y=V(V-1),\qquad N=B^x,
 \qquad q=(BH)^x,\qquad M=q^x.
\]
Initially these power expressions mean requests whose interpretation must
be proved in order. We do not use all these equalities as initial
assumptions. The first one, with nonnegative index \(t\), is available first.
The endpoint identities are
\begin{equation}\label{eq:endpoints}
 x=F+1+(y+1)u_0,\qquad
 x=1+\bigl(1+(t+1)y\bigr)u_t.
\end{equation}
The nine actual CRT representatives are \(a_1,\ldots,a_9\), initially signed.
Replace the old multiplier by
\begin{equation}\label{eq:oddH}
 H=1+2\sum_{j=1}^9a_j.
\end{equation}
This integer is odd, hence nonzero, without any sign information about the
representatives. Once they are nonnegative, it dominates
\(1+\sum a_j\), strengthening the parent's radix estimates.

\subsection{One norm for twelve errors}

The parent already normalizes eleven low binomial remainders to the common
denominator \(M\). Write their signed normalized coordinates as
\(R_1,\ldots,R_{11}\). Their three nontrivial quotient identities include
\[
 q^t v_t=M,\qquad q^{t+1}v_{t+1}=M,\qquad q^{y+1}v_y=M.
\]
The powers here are existing output coordinates, not new exponentiation
operations. The fourth denominator group is \(M\) itself.

Introduce just two new integer coordinates \(d,e\), with quadratic equations
\begin{equation}\label{eq:ceilrows}
 Bd=M,\qquad FM-dV-e=0.
\end{equation}
Use the already present nonnegative universal ternary form
\[
 \G(u,v,w)=u^2+v^2+w(w+1)
\]
to impose the single common bound
\begin{equation}\label{eq:sphere}
 M^2=1+4\left(e^2+\sum_{i=1}^{11}R_i^2\right)+\G(u,v,w).
\end{equation}
The form \(\G\) is nonnegative on integers and represents every nonnegative
integer, as in the retained parent guard construction. A standard
three-square proof is recalled in Section~\ref{sec:complete}.

Delete the three old ternary sign certificates for the ceiling remainder,
the ceiling slack, and the exact-binomial output slack. Delete their now
unused ceiling rows. No sign assertion about \(B,F,q\), or the ceiling error,
is substituted in their place.

Algebraically, if \(s=FB-V\), equations~\eqref{eq:ceilrows} imply
\(e=ds\). Equation~\eqref{eq:sphere} forces \(M\ne0\); therefore
\(B,d\ne0\), and cancellation gives
\begin{equation}\label{eq:rounded-bound}
 |FB-V|=|s|<\frac{|B|}{2}.
\end{equation}
This bound is available before \(B\) has been identified as a binomial
coefficient. It replaces the initial sign assumptions.

\section{A nearest-integer characterization of the factorial}

\begin{theorem}[Rounded factorial identity]\label{thm:rounded}
Let \(t\ge4\) and \(F\in\Z\). Set
\[
 B=\binom{F^2+t}{t},\qquad V=(F^2+1)^t.
\]
Then
\[
 4(FB-V)^2<B^2\quad\Longleftrightarrow\quad F=t!.
\]
Moreover, with \(L=t!\),
\begin{equation}\label{eq:margin}
 0<L-\frac{(L^2+1)^t}{\binom{L^2+t}{t}}
 <\frac{1}{2(t-2)!}.
\end{equation}
For the source range \(t\ge16\), this is in particular less than \(1/32\).
\end{theorem}

\begin{proof}
For \(u\ge0\), put
\[
 f_t(u)=\frac{(u^2+1)^t}{\binom{u^2+t}{t}}
 =L\prod_{i=1}^t\frac{u^2+1}{u^2+i},
 \qquad a=\frac{t(t-1)}2.
\]
Each factor is nondecreasing, \(f_t(0)=1\), and \(0<f_t(u)<L\).
The elementary product inequality \(\prod(1-z_i)\ge1-\sum z_i\),
for \(0\le z_i\le1\), gives
\begin{equation}\label{eq:product-bound}
 f_t(u)\ge L\left(1-\frac{a}{u^2+1}\right).
\end{equation}
The desired strict inequality is exactly \(|F-f_t(|F|)|<1/2\).
It excludes \(F\le0\), and \(f_t(F)<L\) excludes \(F\ge L+1\).

First suppose \(t\ge5\). For \(1\le F\le t-1\), monotonicity gives
\(f_t(F)\ge2^t/(t+1)>t-1/2\ge F+1/2\).
The strict bound holds at \(t=5\) and propagates by induction.
For \(t\le F\le L/2\), equation~\eqref{eq:product-bound} gives
\[
 f_t(F)\ge\frac L2+\frac{L(t+1)}{2(t^2+1)}
 >\frac L2+\frac12\ge F+\frac12.
\]
For \(L/2\le F\le L-1\), it instead gives
\[
 0<L-f_t(F)<\frac{4a}{L}
 =\frac{2}{(t-2)!}\le\frac13<\frac12,
\]
so again \(f_t(F)>F+1/2\). These intervals exclude every positive
integer except \(F=L\). Finally, at \(F=L\), the same product bound gives
\(0<L-f_t(L)<a/L=1/(2(t-2)!)\), proving both the equivalence and the
stronger completeness margin.

It remains to check \(t=4\), where \(L=24\), \(a=6\).
For \(F=1\), \(f_4(1)=16/5>3/2\). For \(F=2,3\), monotonicity and
\(f_4(2)=125/14\) give \(f_4(F)>F+1/2\).
The middle-range argument applies to \(4\le F\le12\).
For \(12\le F\le22\), the error estimate is \(24-f_4(F)<1\),
so \(f_4(F)>23\ge F+1\). For \(F=23\), it improves to
\(24-f_4(23)<144/530<1/2\). The case \(F=24\) has error less than
\(a/L=1/4\). This completes the proof for every \(t\ge4\).
\end{proof}

The uniform threshold cannot be lowered to 2 or 3: at \(t=2\), both
\(F=1\) and \(F=2\) satisfy the rounded inequality; at \(t=3\), both
\(F=5\) and \(F=6\) do. The polynomial construction itself retains
the inherited horizon \(t\ge16\).

The difference between a rounded fixed point and an ordinary ceiling is
material: the new source allows a signed error at the start. Its later
identification with \(t!\) recovers the positive error and its original
interval. The theorem therefore changes the proof order, not merely the
notation of a guard.

\section{The bootstrap without initial sign certificates}\label{sec:signed}

\subsection{The signed-base power lemma}

The inherited congruence has modulus
\(T_n=2An-n^2-1\). Its frame gives
\(|A|>W^W\), \(|n|,|k|,|m|<W\), where \(m\) is the proposed output.
For \(|n|\ge2\), the same proof works with either sign of \(n\):
\(|T_n|>W^W+W\). For \(k\ge0\), the canonical Pell congruence gives
\(m\equiv n^k\pmod{T_n}\); the bounded absolute difference forces equality.
For \(k=-j<0\), it instead gives
\(mn^j\equiv1\pmod{T_n}\), with \(\gcd(n,T_n)=1\) since
\(T_n\equiv-1\pmod n\). The error has absolute value less than
\(W^W+1<|T_n|\), so the bounded difference again forces equality,
which is impossible for integral \(m\) when \(|n|\ge2\) and \(j\ge1\).
Thus \(k\ge0\). The argument uses absolute values throughout. For base 1
and a known nonnegative index, the same congruence gives output 1.

\subsection{Acyclic recovery of magnitudes}

First interpret \(V=(F^2+1)^t\). If \(F=0\), its base is 1 and its index
is the known positive horizon; otherwise its base is at least 2.
Write \(b=|B|\), \(f=|F|\). From~\eqref{eq:rounded-bound},
\begin{equation}\label{eq:basic-magnitude}
 1\le V<(f+1/2)b.
\end{equation}
If \(f=0\), this says \(b>2\). If \(f\ge1\) and \(b\le f\), then
\[
 (f^2+1)^t\ge(f^2+1)^2>f^2+f/2\ge(f+1/2)b,
\]
a contradiction. Consequently
\begin{equation}\label{eq:bootstrap-small}
 |B|\ge2,\qquad |F|<|B|,\qquad V<B^2.
\end{equation}

The request at base \(B\) now forces \(x\ge0\). The odd multiplier
\(H\ne0\) implies \(|BH|\ge2\), so the radix request gives
\(q=(BH)^x\). If \(x=0\), then \(q=1\), and the request for
\(M=q^x\), at base 1 and index 0, gives \(M=1\). This contradicts
\(Bd=M\) and \(|B|\ge2\). Hence \(x\ge1\) and \(|q|\ge2\).

The requests at base \(q\) are now exact. In particular
\(q^{t+1}v_{t+1}=q^x\). The quotient is an integer, so
\begin{equation}\label{eq:xlarge}
 x\ge t+1\ge17,\qquad |q|=|BH|^x\ge|B|^x\ge2^{17}.
\end{equation}
This is the key use of a denominator identity already present in the
normalized construction. It prevents a degenerate bootstrap branch even
when the provisional factorial candidate is zero.

Now \(|q+1|\ge2\). Its power requests force all nine actual CRT
representatives to be nonnegative. Therefore \(H\ge1\), and all remaining
power meanings follow in the same order as in the parent. Since
\(y=V(V-1)\ge0\) and \(x>1\), the second endpoint identity in
\eqref{eq:endpoints} forces \(u_t\ge1\) and
\begin{equation}\label{eq:endpoint-bound}
 x\ge2+(t+1)y.
\end{equation}
At this stage \(B\) and \(q\) may still be negative. No positive-radix
binomial conclusion has yet been assumed.

\section{Signed low digits and the zero factorial branch}

\begin{lemma}[Centered congruence recovery]\label{lem:centered}
Let \(L,M,v\in\Z\) with \(Lv=M\ne0\). If
\(vU=MQ+R\), then \(v\mid R\) and \(r=R/v\) is an integer satisfying
\(U=LQ+r\). Suppose \(|R|<|M|/2\). If another integer
\(r^0\equiv U\pmod L\) satisfies \(|r^0|<|L|/32\), then
\(r=r^0\).
\end{lemma}
\begin{proof}
Divisibility and division recover the old digit equation. The difference
\(r-r^0\) is a multiple of \(L\) with absolute value less than
\((1/2+1/32)|L|<|L|\), hence is zero.
\end{proof}

This lemma does not assume a positive radix or a nonnegative remainder. For
an expansion in a signed base \(q\), the appropriate canonical low value is
the signed polynomial tail, not the least nonnegative remainder.

\subsection{The nonzero factorial candidate}

For integers \(m\ge1\), \(t\ge1\), and \(0\le j\le t\),
\begin{equation}\label{eq:coeffbound}
 \binom{m+t}{j}\le(m+1)^t.
\end{equation}
An induction on \(t\) uses Pascal's identity for the old positions and
\(\binom{m+t+1}{t+1}=\prod_{i=1}^{t+1}(1+m/i)\le(m+1)^{t+1}\)
for the new edge. Apply this with \(m=F^2\ge1\), and put \(Q=|q|\).
The factorial numerator is the exact power \((q+1)^{F^2+t}\), whose low
tail at position \(t\) is
\[
 r^0=\sum_{j=0}^{t-1}\binom{F^2+t}{j}q^j.
\]
Equations~\eqref{eq:bootstrap-small} and~\eqref{eq:xlarge} give
\[
 \frac{|r^0|}{Q^t}<\frac{V}{Q-1}
 <\frac{2}{b^{x-2}}<\frac1{32}.
\]
The common norm and Lemma~\ref{lem:centered} therefore identify the signed
remainder exactly. Dividing its digit equation by \(q^t\) and reducing
modulo \(q\) gives
\[
 B\equiv C:=\binom{F^2+t}{t}\pmod q.
\]
But \(0<C\le V<b^2\), and
\(|B-C|<b+b^2<Q\). Thus \(B=C>0\).

\subsection{The zero factorial candidate needs a different estimate}

One must not use~\eqref{eq:coeffbound} with \(m=0\): it would incorrectly
bound all \(\binom tj\) by 1. Here \(V=1\), \(s=-1\), and the rounded
bound gives \(b\ge3\). Consequently \(Q\ge3^{t+1}\). The valid estimate
\(\binom tj\le2^t\) gives
\[
 \frac{|r^0|}{Q^t}<\frac{2^t}{Q-1}
 <\left(\frac23\right)^{t+1}<\frac1{32}
 \qquad(t\ge16).
\]
The same centered recovery yields \(B\equiv\binom tt=1\pmod q\).
Since \(|B-1|\le b+1<Q\), this forces \(B=1\), contradicting \(b\ge3\).
Thus this branch has no zero.

\subsection{Recovering the factorial and all remaining intervals}

We have proved \(F\ne0\), \(B=\binom{F^2+t}{t}>0\), and
\(|F-V/B|<1/2\). Theorem~\ref{thm:rounded} gives \(F=t!\). In particular
\(F>0\), \(B>F\), \(N=B^x>0\), and \(q=(BH)^x>0\).
The canonical ceiling error is now positive and smaller than \(B/32\).
Hence the deleted ceiling intervals can be restored:
\[
 s=FB-V\ge0,\qquad r=B-s-1\ge0,
 \qquad o=q-B-1\ge0.
\]

For the other ten binomial digits, the parent's coefficient bounds now
apply with the larger multiplier \(H\). Each actual representative is
nonnegative, each required lower index is at most \(x\), and the derived
label satisfies the same degree estimate. The relevant coefficients are
at most \(H^x\), so
\[
 \frac{H^x}{q-1}<\frac2{B^x}<\frac1{32}.
\]
The parent also proves the necessary bound for a zero representative using
\(\max(1,a)\), not \(a\), and proves positivity of the encoded-index
representative before bounding the derived label. These details are
retained unchanged; increasing \(H\) preserves their inequalities.
Lemma~\ref{lem:centered} recovers every genuine low remainder. All range
quotients and arithmetic decoding hypotheses are therefore restored.

\section{Completeness and calibration}\label{sec:complete}

Start with a finite C222 refutation. Use the parent's construction to choose
a sufficiently large horizon \(t=16+h_0^2\), the genuine \(F=t!\), the beta
CRT code \(x\), and the nine actual CRT representatives. Their positivity
and the dominant representative's parity can be chosen as in the parent.
The semantic anchor remains even and dominates all requested indices.

Set \(H=1+2\sum a_i\), \(N=B^x\), \(q=(BH)^x\), and \(M=q^x\).
The larger radix strengthens the radix estimates; it introduces no
new requested exponent. The square frame and private Pell witnesses are
chosen after these finite arithmetic values have been fixed. Thus the
existing completeness construction still applies.

Because \(B\mid q\), the integer \(d=M/B\) exists. Put
\(e=d(FB-V)\). Equation~\eqref{eq:margin} gives
\(0<e<M/32\). The eleven genuine normalized low remainders also lie in
\([0,M/32)\). Therefore
\[
 4\left(e^2+\sum_{i=1}^{11}R_i^2\right)
 <\frac{48}{1024}M^2=\frac3{64}M^2<M^2.
\]
The integer difference
\(M^2-1-4(e^2+\sum R_i^2)\) is nonnegative and has a \(\G\)-representation.
This fills the common norm and gives a zero of the new arithmetic source.

For clarity, universality of \(\G\) follows from the classical three-square
theorem. For \(n\ge0\), \(4n+1\) is neither divisible by 4 nor congruent to
7 modulo 8, so it is a sum of three squares. Exactly one of the three roots
is odd modulo 2; write the other two as \(2u,2v\) and the odd one as
\(2w+1\). Then \(n=u^2+v^2+w(w+1)\). This is the same integer guard
used by the inherited source.

Conversely, Sections~\ref{sec:signed}--\ref{sec:complete} recover the exact
factorial, powers, binomial intervals, and range tests from any new zero.
The inherited amplification estimate and finite Horn-flow decoding then
produce a C222 refutation. The logical calibration is the retained chain
\[
 \text{C222 refutation}\ \longleftrightarrow\
 \text{ZF refutation}\ \longleftrightarrow\
 \text{ZFC refutation},
\]
with the named interfaces
\path{KEndpoint.c222_iff_zf_refutation} and
\path{ConstructibleChoice.zf_refutation_iff_zfc_refutation}.
The interface sources and their pins are retained. The ordinary arithmetic
proofs are formalizable in ZFC and compose with these calibrations to prove
Theorem~\ref{thm:main}.

\paragraph{The precise equivalence being claimed.}
The odd multiplier changes the self-radix equation, so the semantic
transformation is not asserted to be a bijection between all parent and
child zeros with the same front coordinates. Both directions go through
finite refutations. In contrast, constrained monic eliminations have explicit
polynomial restorations and integer coordinate translations are bijective.
The projection also removes 112 unused coordinates, which may be restored
to zero; their removal is not a bijection on all raw assignments.
Keeping these statements separate prevents a false pointwise claim.

\section{Finite polynomial construction and verification}

The raw source edit removes the three initial guard rows, their nine ternary
witnesses, and the obsolete three-row ceiling presentation. It adds the two
quadratic rows~\eqref{eq:ceilrows}, augments the already existing norm by
\(4e^2\), and changes the self-radix multiplier to~\eqref{eq:oddH}.
The useful implementation choice is to encode
\(e=FM-dV\) directly: introducing the old slack \(s\) and then multiplying
\(ds\) would retain a dispensable coordinate.

The source then undergoes the inherited exact Horn factoring and integral
monic projections. Finally, a recorded integer translation chooses a
smaller collected expression. This translation does not change solvability.
The delivered residual list is the final translated one, and is the
authoritative equation in both the appendix and JSON files.

The finite checks cover the following distinct claims:
\begin{itemize}
\item the exact boundary of the source edit, including the removed guards,
the two added rows, the augmented norm, and the odd multiplier;
\item the source-to-factored and factored-to-projected polynomial identities,
including integer restorations of constrained coordinates and zero extensions
for unused coordinates;
\item the final coordinate translation and its inverse;
\item the entire collected expansion of the 181 residual squares, degree,
coefficient height, coordinate liveness, and the printed residual inventory.
\end{itemize}
The independent verifier does not treat a builder's claimed metrics as
evidence. It recomputes the relevant finite identities and expansion.
Normal and optimized Python runs, together with a clean rebuild, are
documented in the delivered receipts. The finite computations do not test
the nonexistence of a zero; the unbounded theorem comes from the proofs.

\paragraph{Reproduction.}
From the package root, run \texttt{python3 verify.py}. Component-specific
commands and the exact rebuild procedure are in \path{README.md} and
\path{direct/README.md}. The polynomial verifier uses the Python standard
library. The PDF is rebuilt by \texttt{sh build\_report.sh} with a standard
TeX Live installation. Appendix generation reads the final JSON list and
introduces no new mathematical coefficients.

\section{A companion equation in ten positive variables}
\label{sec:compactcompanion}

A second construction uses a different tradeoff: ten strictly positive
unknowns, 78 binary arithmetic operations, and exact total degree 454.
It retains the whole four-integer native coefficient tuple of the incoming
\emph{Centered Residues} package. Its matched parent has 76 operations and
degree 464; that package also supplies 77/460. Thus this companion spends
two or one extra operations to reduce degree by ten or six. This is a
specific comparison, not a global simplicity claim. The new argument
recovers the first Pell square only after recovering the main Pell square.

\subsection{The complete polynomial}

The variables are
\[
 J,F,\alpha,i,T,w,s,n,\eta,\zeta\in\mathbb Z_{>0}.
\]
The fixed ports are \(\beta,K_0,M_C,M_F\); their complete literal bytes
and finite reconstruction are supplied in the accompanying compact
component. They obey \(\beta\ge31\) odd, \(K_0>0\), and
\(0<M_C<\beta<M_F<2\beta\). They are constants, not unknowns.
Define the following polynomial abbreviations:
\begin{align*}
 h&=\beta J,&q&=h+1,&C&=q-F-\alpha,\\
 X&=wq,&Y&=q^3s,&k&=\eta+\zeta,\\
 c&=kY+\eta,&A&=Y(X+1)+1,&\Delta&=A^2-1,\\
 H&=4A-5,&\ell&=(XY)(kY),&E&=1+\ell(\ell+k),\\
 R&=[h(q-F)+\alpha](q^2-1)+(M_C+qM_F)J,\\
 e&=k-R-1,& B&=(K_0+w)C-F.
\end{align*}
The new strong-pair subgraph is
\begin{align*}
 D_m&=1+\Delta c^2,&U&=iAc^2,&S&=ED_m,\\
 a&=\Delta U^2S,&F_2&=1+a,&Q&=\Delta a,&G&=SF_2.
\end{align*}
Complete the equation with
\begin{align*}
 b&=cT-R,&V&=F_2b-c,&L&=bV,\\
 K&=X(3c-2X)+2,&M&=hXYH,\\
 Z&=hXY\{H(n-LB)+K\}-H\{XYB+he\},\\
 P&=Q\{Z^2-M^2GV^2\}+(M^2G-Z^2).
\end{align*}
The single equation is \(P=0\). Every abbreviation is polynomial. In
particular \(F_2\) being a square will be a conclusion; its name is not
a square-root instruction. The graph has exactly ten variable leaves
and four fixed coefficient leaves.

\begin{theorem}
Composed with the retained whole-tuple native/refutation calibration,
ordinary mathematics formalizable in ZFC proves
\[
 \bigl(\forall\mathbf x\in\mathbb Z_{>0}^{10},\ P(\mathbf x)\ne0\bigr)
 \quad\Longleftrightarrow\quad\operatorname{Con}(\mathrm{ZFC}).
\]
The polynomial has exact degree 454 and a 78-operation evaluator.
\end{theorem}

\subsection{Delayed recovery of the first square}

For proof notation set \(\delta=Q-1\), \(C_y=QV^2-1\), and
\(z=Z/M\). Positive inputs give \(F_2>2c\), \(Q>F_2>1\), \(M>0\),
and \(|V|>1\), regardless of the initial sign of \(b\). Hence all
radicands in the following argument are positive. The folded identity is
\[
 P=\delta Z^2-M^2GC_y.
\]
Also \(Q=\Delta^2U^2S>c^2\) and \(F_2>S\), estimates used below.
The three required coprimality facts follow directly from the definitions:
\begin{align*}
 &\gcd(S,F_2)=1,
 &&\gcd(S,\delta C_y)=1,\\
 &\delta=\Delta F_2-A^2,\quad F_2\equiv1\pmod A,
 &&\gcd(F_2,\delta)=1,\\
 &C_y\equiv-\Delta c^2-1=-D_m\pmod{F_2},
 \quad F_2\equiv1\pmod{D_m},
 &&\gcd(F_2,C_y)=1.
\end{align*}
The first line uses \(S\mid Q\); the second uses \(A\mid U\); the
last uses \(D_m\mid S\). No coprimality between \(E\) and \(D_m\)
is asserted.

At a zero, \(\delta z^2=GC_y\). Prime valuations and the preceding
coprimalities force \(G\) to be a positive square. Since its positive
factors \(S,F_2\) are coprime, write
\[
 S=\sigma^2,\qquad F_2=f^2,\qquad G=g^2,
 \qquad g=\sigma f>0.
\]
Then \(y=|z|/g\) is a positive rational solution of
\begin{equation}
 QV^2-(Q-1)y^2=1.\label{eq:compactadjacent}
\end{equation}
It is integral. Indeed, if its reduced denominator is \(m\), then
\(m^2\mid Q-1\). For \(v=|V|>1\), one has
\(v<y<Qv/(Q-1)\), and the transformation
\[
 v'=(2Q-1)v-2(Q-1)y,\qquad
 y'=(2Q-1)y-2Qv
\]
preserves the norm and the reduced denominator, has integer \(v'\), and
satisfies \(0<|v'|<v\). Reflect signs and descend to \(|v|=1\), where
\(|y|=1\). Thus \(m=1\), and \(z=\pm gy\) is integral.

Put \(I=U\sigma\). The recovered strong pair satisfies
\[
 f^2-\Delta I^2=1,\qquad Q=(\Delta I)^2,\qquad c^2\mid I.
\]
The retained all-index folding theorem now applies, giving
\(c=\psi_A(p)\) for a positive odd index \(p\). Its hypotheses are the
strong norm, \(0<c\le I\), the integral auxiliary norm
\eqref{eq:compactadjacent}, and \(V\equiv-c\pmod f\). To see the
folding directly, write \(f=\chi_A(m)\), \(I=\psi_A(m)\). Auxiliary
Pell classification yields an odd index \(j\). The standard odd-quotient
identity, reduced modulo \(f\), gives \(\psi_A(j)\equiv\pm c\pmod f\).
Reducing \(j\) modulo \(2m\) and reflecting gives \(0\le p\le m\).
Both \(\psi_A(p)\) and \(c\) are at most \(I<f/2\), so the congruence
cannot wrap and \(c=\psi_A(p)\); parity makes \(p\) odd.

Consequently
\[
 D_m=\chi_A(p)^2=d^2,\qquad E=(\sigma/d)^2.
\]
An integer that is a rational square is an integer square. Hence
\(E=\tau^2\), \(\sigma=\tau d\), and \(g=\tau df\).
This order is essential: the first square is recovered \emph{after}
main folding, not used as an assumption in that folding.

\subsection{Transport, calibration, and the converse}

The first-Pell classification now gives
\(k=2\psi_{2XY^2+1}(n_0)\). Its first index is impossible, because it
would put \(c=2Y+1\) strictly between the first two main ordinates.
Thus
\[
 n_0\ge2,\qquad k\ge8XY^2+4,\qquad c>kY\ge k>8q^7,
 \qquad c>X,
\]
so \(K>0\). The new strong-size estimate is
\[
 \frac{a^2}{g^2}=\Delta U^2\frac a{1+a}
   >\frac{\Delta U^2}{2}>36.
\]
Together with \(a>2c\), it gives \(|V|>3g\) for nonzero \(b\),
\(\operatorname{sgn}V=\operatorname{sgn}b\), and
\(y<3|V|/2\). The case \(b=0\) would make
\((c^2-1)/(Q-1)=y^2-c^2\) an integer strictly between zero and one.

The parent's third-main-index exclusion uses only the unchanged outer
data: if \(p=3\), then \(n_0=2\), and \(0<\eta<k\) forces
\(2Y-X\in\{1,2\}\), contrary to \(q\mid X,Y\). Thus
\(c>4\Delta+3\). In the auxiliary norm, odd indices 1 and 3 are then
excluded. The two signs at index 3 would require
\[
 F_2(b-4\Delta)=c-(4\Delta+3)
 \quad\hbox{or}\quad
 F_2(b+4\Delta)=c+(4\Delta+3),
\]
whose right sides lie strictly between zero and \(2c<F_2\).
Therefore \(|V|>8Q^2\), \(|b|>4F_2>4g\), and
\(gy<3bV/8<L/2\).

For proof notation let
\(\kappa=K/H\), \(\rho=e/(XY)\), \(t=B/h\), and \(W=1+hL\).
The exact carrier is
\[
 z=n+\kappa-Wt-\rho.
\]
Since \(e=b-cT+k-1<b\), one has \(\rho<|b|/(XY)\).
If \(B\le-1\), the carrier and the preceding estimates would give
\(n<L/2-L+|b|/(XY)<0\). Hence \(B\ge0\), which forces \(C>0\).
The inherited mask bound is now licensed: \(3q<R<q^4-2\).
The normalized rank theorem and the signed auxiliary-index theorem
apply to the recovered strong pair, \(c>2p,2R\), and the two residues
\(V\equiv-c\pmod f\), \(V\equiv-R\pmod c\); they yield \(p=R\).

The rest of the retained carrier proof is unchanged. Integrality of \(z\)
gives \(XY\mid hHe\); its divisor \(q^4\), coprime to \(hH\), gives
\(q^4\mid e\). The first recurrence gives
\(k\equiv2n_0\pmod{4XY^2}\). Since \(n_0<R\), the error
\(2n_0-R-1\) is a multiple of \(q^4\) of absolute value below \(q^4\).
Thus \(2n_0=R+1\), and only then the full recurrence proves
\(XY\mid e\). The retained ratio bound and whole-product projection
yield \(X=2^R\), followed by \(h\mid B\). The fixed-minus parity,
half-binomial rounding, and fresh unscaled native reconstruction therefore
have exactly their parent hypotheses and yield a finite ZFC refutation.
The complete inherited proof closure and its coefficient identities are
included alongside \texttt{compact/PROOF.md}.

Conversely, a refutation supplies the parent's canonical positive outer
witnesses, with \(R\equiv3\pmod4\), \(c=\psi_A(R)\),
\(d=\chi_A(R)\), \(E=\tau^2\), and integral
\(\kappa>0,\rho>0,t\ge0\). Choose one positive strong Pell pair
\((f,I)\) with
\[
 Ac^2\tau d\mid I.
\]
It exists by a finite period of the determinant-one Pell matrix modulo
that integer. Set \(i=I/(Ac^2\tau d)>0\). This realizes exactly
\(F_2=f^2\), \(Q=(\Delta I)^2\), and \(G=(\tau df)^2\).
The retained positive auxiliary reconstruction, which needs only
\(c\mid I\), supplies positive \(T,V,y\) for this \emph{same} strong
pair. Finally choose
\[
 n=gy+(1+hL)t+\rho-\kappa.
\]
Here \(\kappa<Xc<f\le g\), so \(n>0\). The carrier equals \(gy\),
proving \(P=0\). Both existential directions are thus established at
the stated inherited calibration boundary.

\subsection{Exact degree and finite verification}

The evaluator has 46 multiplications, 18 additions, and 14 subtractions;
squares and scalar products count. The decisive exact degrees are
\[
 \deg(Q,G,V,M,Z)=(102,134,96,13,117).
\]
Thus \(-QM^2GV^2\) uniquely supplies degree 454, whereas
\(QZ^2\) has degree 336. Writing \(q_0=\beta J\) and
\(k_0=\eta+\zeta\), the entire highest homogeneous part is
\[
 [P]_{454}=-16q_0^{284}w^{42}s^{80}i^8k_0^{38}T^2.
\]
The coefficient of \(J^{284}w^{42}s^{80}i^8\eta^{38}T^2\) is
\(-16\beta^{284}\ne0\) at the actual fixed tuple.
The independent checker covers all 78 gates by 31 exact polynomial cuts,
reconstructs every coefficient byte, checks all 39 terms of this leading
part, and rejects nine mutations that preserve both gate counts and
liveness. Normal and optimized Python runs pass. These finite checks
authenticate the object; the unbounded proof and inherited calibration
remain the mathematical basis of the consistency equivalence.


\section{Further research questions}

\begin{enumerate}
\item \textbf{Can another guard enter the common norm?}
The remaining squared-anchor bound has a different radius and supports the
initial Pell index argument. Sharing its witnesses would require a
noncircular common-denominator theorem. Merely merging inequalities after
they have been interpreted would not justify the earlier anchor step.
\item \textbf{Can the denominator identities do more semantic work?}
The implication \(q^{t+1}\mid q^x\Rightarrow x\ge t+1\) removes a real
bootstrap obstruction. Other divisibility identities may imply bounds now
encoded separately. A useful search should record exactly which power
meanings are already available at each proposed inference.
\item \textbf{Can the rounded factorial certificate simplify the natural
source?}
The signed error is particularly effective over unrestricted integers.
Transferring it to a natural-variable source requires a representation
choice that does not spend more coordinates than it saves.
\item \textbf{Can the two new ceiling coordinates be recycled?}
One represents \(M/B\); the other keeps the scaled ceiling error quadratic.
Removing either without raising degree requires a genuinely shared existing
coordinate or a new arithmetic identity. Direct substitution of the error
into the norm would raise degree.
\item \textbf{Can the new finite transformation be formalized as a local
proof-assistant theorem?}
The rounded factorial lemma, signed-base power lemma, and source-edit
certificate are smaller targets than a fresh end-to-end compiler proof.
They would strengthen the evidence without obscuring the remaining
calibration boundary.
\item \textbf{What lower bounds are appropriate?}
Global simplicity is not one scalar objective. Variable count, degree,
coefficient height, support, arithmetic-circuit length, and the size of the
proof certificate should be reported separately. A lower bound within a
particular encoding class must not be promoted to a lower bound for every
Diophantine presentation of consistency.
\end{enumerate}

\section*{Attribution and evidence status}
\addcontentsline{toc}{section}{Attribution and evidence status}
The common-denominator remainder architecture, ternary guards, signed Pell
orientation, C222 arithmetic, and ZF/ZFC calibration are inherited work,
identified in the retained source ledger. The new material is the joint
ceiling-and-digit norm, the rounded factorial argument, the odd-multiplier
bootstrap, their literal 245-variable realization, and the associated
verification and exposition. For the compact companion, the fixed coefficient
tuple, native calibration, and transport lemmas are inherited; the staged
recovery of the main Pell square and its 78-operation, degree-454 realization
are the new contribution. The research and reviews in this package were
AI-assisted; they are not represented as external human peer review.
No new end-to-end Lean or Rocq theorem is claimed. The comparison is to the
inspected repository snapshot, not a literature-wide priority determination.

\begin{thebibliography}{9}
\bibitem{repo} V. Reshetnikov, \emph{ProveIt}, snapshot
\texttt{3ecd6a6cf53d47519228328567f45fb7ee9781d9}, 25 October 2026.
\url{https://github.com/VladimirReshetnikov/ProveIt/commit/3ecd6a6cf53d47519228328567f45fb7ee9781d9}.
\bibitem{common} \emph{Common remainder certificates for Diophantine
consistency}, incoming research archive, 25 October 2026. Retained here in
\path{direct/inherited/normalized252}; archive and file identities are
recorded in \path{SOURCE_SNAPSHOT.json} and the component source ledger.
\bibitem{forms} \emph{Con(ZFC) in Diophantine form}, collection README in
the snapshot above, \path{Computability/DiophantineConsistency/con-zfc-diophantine-forms}.
\bibitem{sun} Z.-W. Sun, \emph{Ten unknowns for Hilbert's tenth problem
over the integers}, arXiv:2609.01594v2, 2026.
\url{https://arxiv.org/abs/2609.01594}.
\end{thebibliography}

\appendix
\clearpage
\section{Complete literal residual specification}\label{app:residuals}
The variables in this appendix are exactly \(X_0,\ldots,X_{244}\), each
unrestricted over \(\Z\). The polynomial is the sum of the squares of
\emph{all} residuals \(R_0,\ldots,R_{180}\) printed below. A missing monomial
has coefficient zero. The display order is the final JSON order; semantic
register names are not needed to evaluate the polynomial. The companion
\path{variables.tsv} records the construction labels for tracing proofs.
\begingroup
\small
\setlength{\jot}{2pt}
\setlength{\abovedisplayskip}{5pt}
\setlength{\belowdisplayskip}{5pt}
\[\begin{aligned}
R_{0} &={}-X_{3}^{2} + X_{4}
\end{aligned}\]
\[\begin{aligned}
R_{1} &={}X_{0} - X_{0}X_{2} - X_{2} + X_{9}
\end{aligned}\]
\[\begin{aligned}
R_{2} &={}-1 + X_{1} - X_{2}X_{7} - X_{3}
\end{aligned}\]
\[\begin{aligned}
R_{3} &={}-1 + X_{1} - X_{8}X_{9}
\end{aligned}\]
\[\begin{aligned}
R_{4} &={}1 + X_{4} - X_{5}X_{20} + X_{18}
\end{aligned}\]
\[\begin{aligned}
R_{5} &={}-1 + 2X_{10} - X_{10}^{2} - 2X_{10}X_{11} + 2X_{11} - X_{11}^{2} + X_{21}
\end{aligned}\]
\[\begin{aligned}
R_{6} &={}-5 - 30X_{10} - 36X_{10}^{2} - 96X_{10}X_{11} - 24X_{10}X_{21} - 40X_{11}\\
&\quad{} - 64X_{11}^{2} - 32X_{11}X_{21} - 10X_{21} - 4X_{21}^{2} + X_{22}
\end{aligned}\]
\[\begin{aligned}
R_{7} &={}-2 + 2X_{10} - X_{10}^{2} - 4X_{10}X_{22} + 4X_{22} - 4X_{22}^{2} + X_{23}
\end{aligned}\]
\[\begin{aligned}
R_{8} &={}1 + 6X_{10} - 36X_{10}^{2} - 192X_{10}X_{22} - 24X_{10}X_{23} + 16X_{22}\\
&\quad{} - 256X_{22}^{2} - 64X_{22}X_{23} + 2X_{23} - 4X_{23}^{2} + X_{24}
\end{aligned}\]
\[\begin{aligned}
R_{9} &={}-2 - X_{11}^{2} - 2X_{11}X_{12} - X_{12}^{2} + X_{25}
\end{aligned}\]
\[\begin{aligned}
R_{10} &={}-9 + 7X_{10} - X_{10}^{2} - 4X_{10}X_{11} - 4X_{10}X_{25} + 12X_{11}\\
&\quad{} - 4X_{11}^{2} - 8X_{11}X_{25} + 12X_{25} - 4X_{25}^{2} + X_{26}
\end{aligned}\]
\[\begin{aligned}
R_{11} &={}-1 + 2X_{10} - X_{10}^{2} - 2X_{10}X_{12} + 2X_{12} - X_{12}^{2} + X_{27}
\end{aligned}\]
\[\begin{aligned}
R_{12} &={}-6X_{10} - 144X_{10}^{2} - 192X_{10}X_{11} - 192X_{10}X_{12}\\
&\quad{} - 48X_{10}X_{21} - 48X_{10}X_{27} - 8X_{11} - 64X_{11}^{2}\\
&\quad{} - 128X_{11}X_{12} - 32X_{11}X_{21} - 32X_{11}X_{27} - 64X_{12}^{2}\\
&\quad{} - 32X_{12}X_{21} - 32X_{12}X_{27} - 2X_{21} - 4X_{21}^{2} - 8X_{21}X_{27}\\
&\quad{} - 4X_{27}^{2} + X_{28}
\end{aligned}\]
\[\begin{aligned}
R_{13} &={}-4X_{26}^{2} - 8X_{26}X_{28} - 4X_{28}^{2} + X_{29}
\end{aligned}\]
\[\begin{aligned}
R_{14} &={}-6 + 16X_{12} - 16X_{12}^{2} - 16X_{12}X_{25} - 16X_{12}X_{27} + 8X_{25}\\
&\quad{} - 4X_{25}^{2} - 8X_{25}X_{27} + 8X_{27} - 4X_{27}^{2} + X_{30}
\end{aligned}\]
\[\begin{aligned}
R_{15} &={}1 + 6X_{10} - 36X_{10}^{2} - 96X_{10}X_{11} - 48X_{10}X_{12}\\
&\quad{} - 24X_{10}X_{21} - 48X_{10}X_{27} - 24X_{10}X_{30} + 8X_{11} - 64X_{11}^{2}\\
&\quad{} - 64X_{11}X_{12} - 32X_{11}X_{21} - 64X_{11}X_{27} - 32X_{11}X_{30}\\
&\quad{} + 8X_{12} - 16X_{12}^{2} - 16X_{12}X_{21} - 32X_{12}X_{27} - 16X_{12}X_{30}\\
&\quad{} + 2X_{21} - 4X_{21}^{2} - 16X_{21}X_{27} - 8X_{21}X_{30} + 8X_{27}\\
&\quad{} - 16X_{27}^{2} - 16X_{27}X_{30} + 4X_{30} - 4X_{30}^{2} + X_{31}
\end{aligned}\]
\[\begin{aligned}
R_{16} &={}-10 - 5X_{10} - X_{10}^{2} - 12X_{10}X_{11} - 16X_{10}X_{12}\\
&\quad{} - 4X_{10}X_{25} - 36X_{11} - 36X_{11}^{2} - 96X_{11}X_{12} - 24X_{11}X_{25}\\
&\quad{} - 48X_{12} - 64X_{12}^{2} - 32X_{12}X_{25} - 12X_{25} - 4X_{25}^{2} + X_{32}
\end{aligned}\]
\[\begin{aligned}
R_{17} &={}-2 - 24X_{10} - 36X_{10}^{2} - 12X_{10}X_{11} - 96X_{10}X_{12}\\
&\quad{} - 24X_{10}X_{27} - 3X_{11} - X_{11}^{2} - 16X_{11}X_{12} - 4X_{11}X_{27}\\
&\quad{} - 32X_{12} - 64X_{12}^{2} - 32X_{12}X_{27} - 8X_{27} - 4X_{27}^{2} + X_{33}
\end{aligned}\]
\[\begin{aligned}
R_{18} &={}2 - 4X_{32}^{2} - 8X_{32}X_{33} - 4X_{33}^{2} + X_{34}
\end{aligned}\]
\[\begin{aligned}
R_{19} &={}-12 + 16X_{10} - 4X_{10}^{2} - 8X_{10}X_{12} - 8X_{10}X_{21}\\
&\quad{} - 8X_{10}X_{25} + 16X_{12} - 4X_{12}^{2} - 8X_{12}X_{21} - 8X_{12}X_{25}\\
&\quad{} + 16X_{21} - 4X_{21}^{2} - 8X_{21}X_{25} + 16X_{25} - 4X_{25}^{2} + X_{35}
\end{aligned}\]
\[\begin{aligned}
R_{20} &={}-3 - 24X_{10} - 64X_{10}^{2} - 160X_{10}X_{12} - 64X_{10}X_{25}\\
&\quad{} - 32X_{10}X_{27} - 32X_{10}X_{35} - 26X_{12} - 100X_{12}^{2}\\
&\quad{} - 80X_{12}X_{25} - 40X_{12}X_{27} - 40X_{12}X_{35} - 8X_{25} - 16X_{25}^{2}\\
&\quad{} - 16X_{25}X_{27} - 16X_{25}X_{35} - 6X_{27} - 4X_{27}^{2} - 8X_{27}X_{35}\\
&\quad{} - 4X_{35} - 4X_{35}^{2} + X_{36}
\end{aligned}\]
\[\begin{aligned}
R_{21} &={}-1 - 18X_{10} - 36X_{10}^{2} - 96X_{10}X_{11} - 24X_{10}X_{21} - 24X_{11}\\
&\quad{} - 64X_{11}^{2} - 32X_{11}X_{21} - 6X_{21} - 4X_{21}^{2} + X_{37}
\end{aligned}\]
\[\begin{aligned}
R_{22} &={}-2 + 2X_{10} - X_{10}^{2} - 4X_{10}X_{37} + 4X_{37} - 4X_{37}^{2} + X_{38}
\end{aligned}\]
\[\begin{aligned}
R_{23} &={}1 + 6X_{10} - 36X_{10}^{2} - 192X_{10}X_{37} - 24X_{10}X_{38} + 16X_{37}\\
&\quad{} - 256X_{37}^{2} - 64X_{37}X_{38} + 2X_{38} - 4X_{38}^{2} + X_{39}
\end{aligned}\]
\[\begin{aligned}
R_{24} &={}-X_{11}^{2} - 4X_{11}X_{39} - 4X_{39}^{2} + X_{40}
\end{aligned}\]
\[\begin{aligned}
R_{25} &={}-6 - 30X_{10} - 64X_{10}^{2} - 128X_{10}X_{11} - 64X_{10}X_{21} - 32X_{11}\\
&\quad{} - 64X_{11}^{2} - 64X_{11}X_{21} - 14X_{21} - 16X_{21}^{2} + X_{41}
\end{aligned}\]
\[\begin{aligned}
R_{26} &={}-4 + 8X_{11} - 64X_{11}^{2} - 96X_{11}X_{12} - 32X_{11}X_{25}\\
&\quad{} - 32X_{11}X_{41} + 6X_{12} - 36X_{12}^{2} - 24X_{12}X_{25} - 24X_{12}X_{41}\\
&\quad{} + 2X_{25} - 4X_{25}^{2} - 8X_{25}X_{41} + 4X_{41} - 4X_{41}^{2} + X_{42}
\end{aligned}\]
\[\begin{aligned}
R_{27} &={}8X_{10} - 64X_{10}^{2} - 96X_{10}X_{12} - 32X_{10}X_{27} - 32X_{10}X_{42}\\
&\quad{} + 6X_{12} - 36X_{12}^{2} - 24X_{12}X_{27} - 24X_{12}X_{42} + 2X_{27}\\
&\quad{} - 4X_{27}^{2} - 8X_{27}X_{42} + 4X_{42} - 4X_{42}^{2} + X_{43}
\end{aligned}\]
\[\begin{aligned}
R_{28} &={}-X_{12}^{2} - 2X_{12}X_{13} - X_{13}^{2} + X_{44}
\end{aligned}\]
\[\begin{aligned}
R_{29} &={}2X_{10} - X_{10}^{2} - 4X_{10}X_{12} - 4X_{10}X_{44} + 4X_{12} - 4X_{12}^{2}\\
&\quad{} - 8X_{12}X_{44} + 4X_{44} - 4X_{44}^{2} + X_{45}
\end{aligned}\]
\[\begin{aligned}
R_{30} &={}-35 + 24X_{10} - 4X_{10}^{2} - 8X_{10}X_{11} - 8X_{10}X_{25}\\
&\quad{} - 8X_{10}X_{27} + 24X_{11} - 4X_{11}^{2} - 8X_{11}X_{25} - 8X_{11}X_{27}\\
&\quad{} + 24X_{25} - 4X_{25}^{2} - 8X_{25}X_{27} + 24X_{27} - 4X_{27}^{2} + X_{46}
\end{aligned}\]
\[\begin{aligned}
R_{31} &={}-16X_{10}^{2} - 32X_{10}X_{12} - 32X_{10}X_{27} - 32X_{10}X_{44}\\
&\quad{} - 16X_{10}X_{45} - 16X_{10}X_{46} - 4X_{12} - 16X_{12}^{2} - 32X_{12}X_{27}\\
&\quad{} - 32X_{12}X_{44} - 16X_{12}X_{45} - 16X_{12}X_{46} - 16X_{27}^{2}\\
&\quad{} - 32X_{27}X_{44} - 16X_{27}X_{45} - 16X_{27}X_{46} - 4X_{44} - 16X_{44}^{2}\\
&\quad{} - 16X_{44}X_{45} - 16X_{44}X_{46} - 2X_{45} - 4X_{45}^{2} - 8X_{45}X_{46}\\
&\quad{} - 4X_{46}^{2} + X_{47}
\end{aligned}\]
\[\begin{aligned}
R_{32} &={}X_{12} - 25X_{12}^{2} - 40X_{12}X_{27} - 20X_{12}X_{30} - 16X_{27}^{2}\\
&\quad{} - 16X_{27}X_{30} - 4X_{30}^{2} + X_{48}
\end{aligned}\]
\[\begin{aligned}
R_{33} &={}-1 + 8X_{11} - 4X_{11}^{2} - 8X_{11}X_{25} + 8X_{25} - 4X_{25}^{2} + X_{49}
\end{aligned}\]
\[\begin{aligned}
R_{34} &={}-1 - 16X_{11}^{2} - 8X_{11}X_{12} - 32X_{11}X_{25} - 16X_{11}X_{49}\\
&\quad{} - X_{12}^{2} - 8X_{12}X_{25} - 4X_{12}X_{49} - 16X_{25}^{2} - 16X_{25}X_{49}\\
&\quad{} - 4X_{49}^{2} + X_{50}
\end{aligned}\]
\[\begin{aligned}
R_{35} &={}1 + X_{11} - X_{11}^{2} - 4X_{11}X_{12} - 4X_{11}X_{50} - 4X_{12}^{2}\\
&\quad{} - 8X_{12}X_{50} - 4X_{50}^{2} + X_{51}
\end{aligned}\]
\[\begin{aligned}
R_{36} &={}-1 - 4X_{51} - 4X_{51}^{2} + X_{52}
\end{aligned}\]
\[\begin{aligned}
R_{37} &={}-1 - X_{12}^{2} - 8X_{12}X_{51} - 4X_{12}X_{52} - 16X_{51}^{2}\\
&\quad{} - 16X_{51}X_{52} - 4X_{52}^{2} + X_{53}
\end{aligned}\]
\[\begin{aligned}
R_{38} &={}2 + X_{11} - X_{11}^{2} - 4X_{11}X_{12} - 4X_{11}X_{53} - 4X_{12}^{2}\\
&\quad{} - 8X_{12}X_{53} - 4X_{53}^{2} + X_{54}
\end{aligned}\]
\[\begin{aligned}
R_{39} &={}-4X_{48}^{2} - 8X_{48}X_{54} - 4X_{54}^{2} + X_{55}
\end{aligned}\]
\[\begin{aligned}
R_{40} &={}-9 + 7X_{10} - X_{10}^{2} - 8X_{10}X_{48} - 4X_{10}X_{55} + 24X_{48}\\
&\quad{} - 16X_{48}^{2} - 16X_{48}X_{55} + 12X_{55} - 4X_{55}^{2} + X_{56}
\end{aligned}\]
\[\begin{aligned}
R_{41} &={}-1 - 4X_{56} - 4X_{56}^{2} + X_{57}
\end{aligned}\]
\[\begin{aligned}
R_{42} &={}-X_{12}^{2} - 8X_{12}X_{56} - 4X_{12}X_{57} - 16X_{56}^{2} - 16X_{56}X_{57}\\
&\quad{} - 4X_{57}^{2} + X_{58}
\end{aligned}\]
\[\begin{aligned}
R_{43} &={}-X_{10} - X_{10}^{2} - 4X_{10}X_{12} - 4X_{10}X_{58} - 4X_{12} - 4X_{12}^{2}\\
&\quad{} - 8X_{12}X_{58} - 4X_{58} - 4X_{58}^{2} + X_{59}
\end{aligned}\]
\[\begin{aligned}
R_{44} &={}-10 - 12X_{59} - 4X_{59}^{2} + X_{60}
\end{aligned}\]
\[\begin{aligned}
R_{45} &={}1 + 4X_{11} - 4X_{11}^{2} - 8X_{11}X_{25} - 16X_{11}X_{59}\\
&\quad{} - 8X_{11}X_{60} + 4X_{25} - 4X_{25}^{2} - 16X_{25}X_{59} - 8X_{25}X_{60}\\
&\quad{} + 8X_{59} - 16X_{59}^{2} - 16X_{59}X_{60} + 4X_{60} - 4X_{60}^{2} + X_{61}
\end{aligned}\]
\[\begin{aligned}
R_{46} &={}X_{11} - 25X_{11}^{2} - 40X_{11}X_{25} - 20X_{11}X_{61} - 16X_{25}^{2}\\
&\quad{} - 16X_{25}X_{61} - 4X_{61}^{2} + X_{62}
\end{aligned}\]
\[\begin{aligned}
R_{47} &={}-10 - 12X_{54} - 4X_{54}^{2} + X_{63}
\end{aligned}\]
\[\begin{aligned}
R_{48} &={}1 - 16X_{54}^{2} - 16X_{54}X_{62} - 16X_{54}X_{63} - 4X_{62}^{2}\\
&\quad{} - 8X_{62}X_{63} - 4X_{63}^{2} + X_{64}
\end{aligned}\]
\[\begin{aligned}
R_{49} &={}X_{12} - X_{12}^{2} - 8X_{12}X_{62} - 4X_{12}X_{64} - 16X_{62}^{2}\\
&\quad{} - 16X_{62}X_{64} - 4X_{64}^{2} + X_{65}
\end{aligned}\]
\[\begin{aligned}
R_{50} &={}-4X_{10} - 4X_{10}^{2} - 4X_{10}X_{12} - 8X_{10}X_{21} - X_{12} - X_{12}^{2}\\
&\quad{} - 4X_{12}X_{21} - 4X_{21} - 4X_{21}^{2} + X_{66}
\end{aligned}\]
\[\begin{aligned}
R_{51} &={}1 + 4X_{10} - 4X_{10}^{2} - 8X_{10}X_{21} - 8X_{10}X_{66} + 4X_{21}\\
&\quad{} - 4X_{21}^{2} - 8X_{21}X_{66} + 4X_{66} - 4X_{66}^{2} + X_{67}
\end{aligned}\]
\[\begin{aligned}
R_{52} &={}-2 - 4X_{10} - 16X_{10}^{2} - 32X_{10}X_{21} - 16X_{10}X_{65}\\
&\quad{} - 16X_{10}X_{67} - 4X_{21} - 16X_{21}^{2} - 16X_{21}X_{65} - 16X_{21}X_{67}\\
&\quad{} - 4X_{65}^{2} - 8X_{65}X_{67} - 2X_{67} - 4X_{67}^{2} + X_{68}
\end{aligned}\]
\[\begin{aligned}
R_{53} &={}-2 - 4X_{68} - 4X_{68}^{2} + X_{69}
\end{aligned}\]
\[\begin{aligned}
R_{54} &={}-X_{3}X_{15} + X_{70}
\end{aligned}\]
\[\begin{aligned}
R_{55} &={}X_{14} - X_{15} - 2X_{24}X_{70} - X_{71}
\end{aligned}\]
\[\begin{aligned}
R_{56} &={}X_{14} - X_{15} - 4X_{26}X_{70} - 2X_{29}X_{70} - X_{72}
\end{aligned}\]
\[\begin{aligned}
R_{57} &={}X_{14} - X_{15} - 2X_{31}X_{70} - 6X_{70} - X_{73}
\end{aligned}\]
\[\begin{aligned}
R_{58} &={}X_{14} - X_{15} - 4X_{32}X_{70} - 2X_{34}X_{70} - X_{74}
\end{aligned}\]
\[\begin{aligned}
R_{59} &={}X_{14} - X_{15} - 2X_{36}X_{70} + 6X_{70} - X_{75}
\end{aligned}\]
\[\begin{aligned}
R_{60} &={}X_{14} - X_{15} - 4X_{39}X_{70} - 2X_{40}X_{70} - X_{76}
\end{aligned}\]
\[\begin{aligned}
R_{61} &={}X_{14} - X_{15} - 2X_{43}X_{70} - X_{77}
\end{aligned}\]
\[\begin{aligned}
R_{62} &={}X_{14} - X_{15} - 2X_{47}X_{70} - 4X_{70} - X_{78}
\end{aligned}\]
\[\begin{aligned}
R_{63} &={}-X_{10}X_{11} + X_{10}X_{12} + X_{11} + X_{11}X_{12} - X_{12}\\
&\quad{} - X_{12}^{2} + X_{79}
\end{aligned}\]
\[\begin{aligned}
R_{64} &={}1 + 2X_{13} + X_{13}^{2} - X_{79}^{2} + X_{80}
\end{aligned}\]
\[\begin{aligned}
R_{65} &={}-X_{80}^{2} + X_{81}
\end{aligned}\]
\[\begin{aligned}
R_{66} &={}-4X_{68}X_{70} - 2X_{69}X_{70} + 2X_{70} + X_{82}
\end{aligned}\]
\[\begin{aligned}
R_{67} &={}-X_{14}^{2} + 2X_{14}X_{15} + 2X_{14}X_{82} - X_{15}^{2} - 2X_{15}X_{82}\\
&\quad{} - X_{82}^{2} + X_{83}
\end{aligned}\]
\[\begin{aligned}
R_{68} &={}-1 - X_{3}X_{11} + X_{84}
\end{aligned}\]
\[\begin{aligned}
R_{69} &={}-X_{15}X_{84} + X_{86}X_{87} + X_{88}
\end{aligned}\]
\[\begin{aligned}
R_{70} &={}-1 + X_{3} - X_{3}X_{10} + X_{85}
\end{aligned}\]
\[\begin{aligned}
R_{71} &={}-1 + 2X_{3} - 2X_{3}X_{10} - 2X_{3}X_{21} + X_{86}
\end{aligned}\]
\[\begin{aligned}
R_{72} &={}-X_{14}X_{85} + X_{87}
\end{aligned}\]
\[\begin{aligned}
R_{73} &={}-1 - 6X_{3}X_{10} - 8X_{3}X_{11} - 2X_{3}X_{21} + X_{89}
\end{aligned}\]
\[\begin{aligned}
R_{74} &={}X_{14}X_{84} - X_{15}X_{89} + X_{90}
\end{aligned}\]
\[\begin{aligned}
R_{75} &={}-X_{14}X_{71} + X_{15}X_{71} + X_{91}
\end{aligned}\]
\[\begin{aligned}
R_{76} &={}-X_{72}X_{91} + X_{92}
\end{aligned}\]
\[\begin{aligned}
R_{77} &={}-X_{73}X_{92} + X_{93}
\end{aligned}\]
\[\begin{aligned}
R_{78} &={}-X_{74}X_{93} + X_{94}
\end{aligned}\]
\[\begin{aligned}
R_{79} &={}-X_{75}X_{94} + X_{95}
\end{aligned}\]
\[\begin{aligned}
R_{80} &={}-X_{76}X_{95} + X_{96}
\end{aligned}\]
\[\begin{aligned}
R_{81} &={}-X_{77}X_{96} + X_{97}
\end{aligned}\]
\[\begin{aligned}
R_{82} &={}-X_{78}X_{97} + X_{98}
\end{aligned}\]
\[\begin{aligned}
R_{83} &={}-X_{88}X_{98} + X_{99}
\end{aligned}\]
\[\begin{aligned}
R_{84} &={}-X_{90}X_{99} + X_{100}
\end{aligned}\]
\[\begin{aligned}
R_{85} &={}-X_{100}^{2} + X_{101}
\end{aligned}\]
\[\begin{aligned}
R_{86} &={}-X_{2} - X_{2}X_{18} + X_{18} + X_{102}
\end{aligned}\]
\[\begin{aligned}
R_{87} &={}-X_{16}X_{102} + X_{103}
\end{aligned}\]
\[\begin{aligned}
R_{88} &={}-2X_{1}^{2} + 2X_{1}X_{14} + 2X_{1}X_{15} + 2X_{1}X_{103}\\
&\quad{} + 2X_{1}X_{104} - X_{14}^{2} - 2X_{14}X_{103} - X_{15}^{2} - 2X_{15}X_{104}\\
&\quad{} + X_{19}X_{106} - X_{81}X_{101} - X_{83}X_{101} - X_{103}^{2} - X_{104}^{2}
\end{aligned}\]
\[\begin{aligned}
R_{89} &={}-X_{2}X_{17} + X_{17} - X_{17}X_{102} + X_{104}
\end{aligned}\]
\[\begin{aligned}
R_{90} &={}-9X_{10} - 9X_{11} - 4X_{21} - 2X_{24} - 4X_{26} - 2X_{29} - 2X_{31}\\
&\quad{} - 4X_{32} - 2X_{34} - 2X_{36} - 4X_{39} - 2X_{40} - 2X_{43} - 2X_{47}\\
&\quad{} - 4X_{68} - 2X_{69} + X_{105}
\end{aligned}\]
\[\begin{aligned}
R_{91} &={}-X_{107}X_{110} + X_{110} + X_{112}
\end{aligned}\]
\[\begin{aligned}
R_{92} &={}-X_{5}X_{134} + X_{108}X_{234} - X_{111} - X_{112}X_{134}
\end{aligned}\]
\[\begin{aligned}
R_{93} &={}X_{113}X_{235} - X_{115}X_{137} - X_{116}
\end{aligned}\]
\[\begin{aligned}
R_{94} &={}X_{117}X_{235} - X_{118}X_{137} - X_{119}
\end{aligned}\]
\[\begin{aligned}
R_{95} &={}X_{120}X_{235} - X_{121}X_{137} - X_{122}
\end{aligned}\]
\[\begin{aligned}
R_{96} &={}X_{123}X_{235} - X_{124}X_{137} - X_{125}
\end{aligned}\]
\[\begin{aligned}
R_{97} &={}X_{126}X_{236} - X_{128}X_{137} - X_{129}
\end{aligned}\]
\[\begin{aligned}
R_{98} &={}X_{130}X_{236} - X_{131}X_{137} - X_{132}
\end{aligned}\]
\[\begin{aligned}
R_{99} &={}-X_{19}X_{134} + X_{137}
\end{aligned}\]
\[\begin{aligned}
R_{100} &={}X_{133} - X_{135}X_{137} - X_{136}
\end{aligned}\]
\[\begin{aligned}
R_{101} &={}-X_{137}X_{139} + X_{138} - X_{140}
\end{aligned}\]
\[\begin{aligned}
R_{102} &={}-X_{137}X_{142} + X_{141}X_{234} - X_{143}
\end{aligned}\]
\[\begin{aligned}
R_{103} &={}-X_{137}X_{145} + X_{144}X_{235} - X_{146}
\end{aligned}\]
\[\begin{aligned}
R_{104} &={}X_{147} - 2X_{147}X_{149} + X_{150}
\end{aligned}\]
\[\begin{aligned}
R_{105} &={}-X_{147}X_{150} - X_{150} + X_{151}
\end{aligned}\]
\[\begin{aligned}
R_{106} &={}2X_{148} + X_{148}^{2} + X_{150}^{2} - X_{151}^{2}
\end{aligned}\]
\[\begin{aligned}
R_{107} &={}-X_{148}X_{153} - X_{153} + X_{154}
\end{aligned}\]
\[\begin{aligned}
R_{108} &={}2X_{152} + X_{152}^{2} + X_{153}^{2} - X_{154}^{2}
\end{aligned}\]
\[\begin{aligned}
R_{109} &={}X_{4}^{2} - 2X_{4}X_{148} - 2X_{148} + X_{155}
\end{aligned}\]
\[\begin{aligned}
R_{110} &={}-1 - X_{4}X_{153} + X_{6} - X_{152} - X_{153} + X_{154}\\
&\quad{} + X_{155}X_{156}
\end{aligned}\]
\[\begin{aligned}
R_{111} &={}1 - 2X_{5} + X_{5}^{2} - 2X_{5}X_{148} + X_{157}
\end{aligned}\]
\[\begin{aligned}
R_{112} &={}-X_{5}X_{203} + X_{19} + X_{157}X_{158} - X_{202} + X_{204}
\end{aligned}\]
\[\begin{aligned}
R_{113} &={}-1 + 2X_{148}X_{218} - X_{159} + 2X_{218} - X_{218}^{2}
\end{aligned}\]
\[\begin{aligned}
R_{114} &={}-1 + X_{107} + X_{159}X_{160} - X_{202} - X_{203}X_{218} + X_{204}
\end{aligned}\]
\[\begin{aligned}
R_{115} &={}-X_{148}X_{162} - X_{162} + X_{163}
\end{aligned}\]
\[\begin{aligned}
R_{116} &={}2X_{161} + X_{161}^{2} + X_{162}^{2} - X_{163}^{2}
\end{aligned}\]
\[\begin{aligned}
R_{117} &={}1 - 2X_{107} + X_{107}^{2} - 2X_{107}X_{148} + X_{164}
\end{aligned}\]
\[\begin{aligned}
R_{118} &={}-1 - X_{107}X_{162} + X_{108} - X_{161} + X_{163} + X_{164}X_{165}
\end{aligned}\]
\[\begin{aligned}
R_{119} &={}4 - 4X_{107} + X_{107}^{2} - 2X_{107}X_{148} + 2X_{148} + X_{166}
\end{aligned}\]
\[\begin{aligned}
R_{120} &={}-1 - X_{107}X_{153} + X_{109} - X_{152} + X_{153} + X_{154}\\
&\quad{} + X_{166}X_{167}
\end{aligned}\]
\[\begin{aligned}
R_{121} &={}1 + X_{148} - X_{148}X_{169} - X_{169} + X_{170}
\end{aligned}\]
\[\begin{aligned}
R_{122} &={}1 + 2X_{168} + X_{168}^{2} - 2X_{169} + X_{169}^{2} - X_{170}^{2}
\end{aligned}\]
\[\begin{aligned}
R_{123} &={}-1 + X_{107} - X_{107}X_{169} + X_{113} + X_{164}X_{171} - X_{168}\\
&\quad{} + X_{170}
\end{aligned}\]
\[\begin{aligned}
R_{124} &={}-2X_{148} - X_{148}^{2} + X_{172}
\end{aligned}\]
\[\begin{aligned}
R_{125} &={}1 - X_{107} - X_{107}X_{152} - X_{107}X_{154} + X_{114} + X_{148}X_{154}\\
&\quad{} + X_{152} - X_{153}X_{172} + 2X_{154} + X_{166}X_{173}
\end{aligned}\]
\[\begin{aligned}
R_{126} &={}-X_{148}X_{175} - X_{175} + X_{176}
\end{aligned}\]
\[\begin{aligned}
R_{127} &={}2X_{174} + X_{174}^{2} + X_{175}^{2} - X_{176}^{2}
\end{aligned}\]
\[\begin{aligned}
R_{128} &={}-1 - X_{107}X_{175} + X_{117} + X_{164}X_{177} - X_{174} + X_{176}
\end{aligned}\]
\[\begin{aligned}
R_{129} &={}-X_{148}X_{179} - X_{179} + X_{180}
\end{aligned}\]
\[\begin{aligned}
R_{130} &={}2X_{178} + X_{178}^{2} + X_{179}^{2} - X_{180}^{2}
\end{aligned}\]
\[\begin{aligned}
R_{131} &={}-1 - X_{107}X_{179} + X_{120} + X_{164}X_{181} - X_{178} + X_{180}
\end{aligned}\]
\[\begin{aligned}
R_{132} &={}-X_{148}X_{183} - X_{183} + X_{184}
\end{aligned}\]
\[\begin{aligned}
R_{133} &={}2X_{182} + X_{182}^{2} + X_{183}^{2} - X_{184}^{2}
\end{aligned}\]
\[\begin{aligned}
R_{134} &={}-1 - X_{107}X_{183} + X_{123} + X_{164}X_{185} - X_{182} + X_{184}
\end{aligned}\]
\[\begin{aligned}
R_{135} &={}-X_{148}X_{187} - X_{187} + X_{188}
\end{aligned}\]
\[\begin{aligned}
R_{136} &={}2X_{186} + X_{186}^{2} + X_{187}^{2} - X_{188}^{2}
\end{aligned}\]
\[\begin{aligned}
R_{137} &={}-1 - X_{107}X_{187} + X_{126} + X_{164}X_{189} - X_{186} + X_{188}
\end{aligned}\]
\[\begin{aligned}
R_{138} &={}-X_{148}X_{191} - X_{191} + X_{192}
\end{aligned}\]
\[\begin{aligned}
R_{139} &={}2X_{190} + X_{190}^{2} + X_{191}^{2} - X_{192}^{2}
\end{aligned}\]
\[\begin{aligned}
R_{140} &={}-1 - X_{107}X_{191} + X_{127} + X_{166}X_{193} - X_{190} + X_{191}\\
&\quad{} + X_{192}
\end{aligned}\]
\[\begin{aligned}
R_{141} &={}-X_{148}X_{195} - X_{195} + X_{196}
\end{aligned}\]
\[\begin{aligned}
R_{142} &={}2X_{194} + X_{194}^{2} + X_{195}^{2} - X_{196}^{2}
\end{aligned}\]
\[\begin{aligned}
R_{143} &={}-1 - X_{107}X_{195} + X_{130} + X_{164}X_{197} - X_{194} + X_{196}
\end{aligned}\]
\[\begin{aligned}
R_{144} &={}-X_{148}X_{199} - X_{199} + X_{200}
\end{aligned}\]
\[\begin{aligned}
R_{145} &={}2X_{198} + X_{198}^{2} + X_{199}^{2} - X_{200}^{2}
\end{aligned}\]
\[\begin{aligned}
R_{146} &={}-1 - X_{107}X_{199} + X_{133} + X_{164}X_{201} - X_{198} + X_{200}
\end{aligned}\]
\[\begin{aligned}
R_{147} &={}1 - X_{148}X_{203} - X_{203} + X_{204}
\end{aligned}\]
\[\begin{aligned}
R_{148} &={}-1 + 2X_{202} + X_{202}^{2} + X_{203}^{2} - 2X_{204} - X_{204}^{2}
\end{aligned}\]
\[\begin{aligned}
R_{149} &={}-X_{107}X_{203} + X_{134} + X_{166}X_{205} - X_{202} + X_{203} + X_{204}
\end{aligned}\]
\[\begin{aligned}
R_{150} &={}-X_{148}X_{207} - X_{207} + X_{208}
\end{aligned}\]
\[\begin{aligned}
R_{151} &={}2X_{206} + X_{206}^{2} + X_{207}^{2} - X_{208}^{2}
\end{aligned}\]
\[\begin{aligned}
R_{152} &={}-1 - X_{107}X_{207} + X_{138} + X_{164}X_{209} - X_{206} + X_{208}
\end{aligned}\]
\[\begin{aligned}
R_{153} &={}-X_{148}X_{211} - X_{211} + X_{212}
\end{aligned}\]
\[\begin{aligned}
R_{154} &={}2X_{210} + X_{210}^{2} + X_{211}^{2} - X_{212}^{2}
\end{aligned}\]
\[\begin{aligned}
R_{155} &={}-1 - X_{107}X_{211} + X_{141} + X_{164}X_{213} - X_{210} + X_{212}
\end{aligned}\]
\[\begin{aligned}
R_{156} &={}-X_{148}X_{215} - X_{215} + X_{216}
\end{aligned}\]
\[\begin{aligned}
R_{157} &={}2X_{214} + X_{214}^{2} + X_{215}^{2} - X_{216}^{2}
\end{aligned}\]
\[\begin{aligned}
R_{158} &={}-1 - X_{107}X_{215} + X_{144} + X_{164}X_{217} - X_{214} + X_{216}
\end{aligned}\]
\[\begin{aligned}
R_{159} &={}-6 - 2X_{0} - 3X_{0}^{2} - 2X_{0}X_{4} - X_{1}^{2} - X_{2}^{2} - 2X_{4}\\
&\quad{} - 2X_{4}^{2} - X_{5}^{2} - X_{6}^{2} + 2X_{10} - X_{10}^{2} - X_{11}^{2} - X_{12}^{2}\\
&\quad{} - X_{13}^{2} - X_{14}^{2} - X_{15}^{2} - X_{16}^{2} - X_{17}^{2} - X_{18}^{2} - X_{19}^{2}\\
&\quad{} - X_{105}^{2} + 2X_{107} - 2X_{107}^{2} - X_{108}^{2} - X_{109}^{2} - X_{113}^{2} - X_{114}^{2}\\
&\quad{} - X_{117}^{2} - X_{120}^{2} - X_{123}^{2} - X_{126}^{2} - X_{127}^{2} - X_{130}^{2} - X_{133}^{2}\\
&\quad{} - X_{134}^{2} - X_{138}^{2} - X_{141}^{2} - X_{144}^{2} + X_{147} - X_{218}^{2}
\end{aligned}\]
\[\begin{aligned}
R_{160} &={}X_{5} - 2X_{5}X_{10} - 2X_{5}X_{11} - 2X_{5}X_{12}\\
&\quad{} - 2X_{5}X_{13} - 2X_{5}X_{14} - 2X_{5}X_{15} - 2X_{5}X_{16}\\
&\quad{} - 2X_{5}X_{17} - 2X_{5}X_{18} + X_{218}
\end{aligned}\]
\[\begin{aligned}
R_{161} &={}-16 + X_{0} - X_{219}^{2}
\end{aligned}\]
\[\begin{aligned}
R_{162} &={}-1 + X_{2} + X_{6} - X_{6}^{2}
\end{aligned}\]
\[\begin{aligned}
R_{163} &={}-X_{148}X_{221} - X_{221} + X_{230}
\end{aligned}\]
\[\begin{aligned}
R_{164} &={}-1 + X_{220}^{2} + X_{221}^{2} - X_{230}^{2}
\end{aligned}\]
\[\begin{aligned}
R_{165} &={}-X_{223} - X_{223}X_{224} + X_{231}
\end{aligned}\]
\[\begin{aligned}
R_{166} &={}-1 + X_{222}^{2} + X_{223}^{2} - X_{231}^{2}
\end{aligned}\]
\[\begin{aligned}
R_{167} &={}-X_{215}^{2} + X_{232}
\end{aligned}\]
\[\begin{aligned}
R_{168} &={}X_{221} - 2X_{226}X_{232} - X_{232}
\end{aligned}\]
\[\begin{aligned}
R_{169} &={}-4X_{221} - 8X_{221}X_{225} + X_{224}
\end{aligned}\]
\[\begin{aligned}
R_{170} &={}-X_{148} - X_{220}X_{227} + X_{224}
\end{aligned}\]
\[\begin{aligned}
R_{171} &={}-1 - X_{214} - X_{220}X_{228} + X_{222}
\end{aligned}\]
\[\begin{aligned}
R_{172} &={}-X_{105} - 4X_{215}X_{229} + X_{223}
\end{aligned}\]
\[\begin{aligned}
R_{173} &={}2X_{0}^{2} + 2X_{0}X_{4} - 2X_{0}X_{153} - 2X_{0}X_{162} + X_{1}^{2}\\
&\quad{} - 2X_{1}X_{203} + X_{2}^{2} - 2X_{2}X_{191} + X_{4}^{2} - 2X_{4}X_{162}\\
&\quad{} + X_{10}^{2} - 2X_{10}X_{169} + X_{11}^{2} - 2X_{11}X_{175} + X_{12}^{2}\\
&\quad{} - 2X_{12}X_{179} + X_{13}^{2} - 2X_{13}X_{183} + X_{14}^{2} - 2X_{14}X_{187}\\
&\quad{} + X_{15}^{2} - 2X_{15}X_{195} + X_{16}^{2} - 2X_{16}X_{199} + X_{17}^{2}\\
&\quad{} - 2X_{17}X_{207} + X_{18}^{2} - 2X_{18}X_{211} + X_{105}^{2} - 2X_{105}X_{215}\\
&\quad{} - X_{148}X_{233} + X_{152}^{2} + X_{153}^{2} + X_{161}^{2} + X_{162}^{2} + X_{168}^{2}\\
&\quad{} + X_{169}^{2} + X_{174}^{2} + X_{175}^{2} + X_{178}^{2} + X_{179}^{2} + X_{182}^{2} + X_{183}^{2}\\
&\quad{} + X_{186}^{2} + X_{187}^{2} + X_{190}^{2} + X_{191}^{2} + X_{194}^{2} + X_{195}^{2} + X_{198}^{2}\\
&\quad{} + X_{199}^{2} + X_{202}^{2} + X_{203}^{2} + X_{206}^{2} + X_{207}^{2} + X_{210}^{2} + X_{211}^{2}\\
&\quad{} + X_{214}^{2} + X_{215}^{2}
\end{aligned}\]
\[\begin{aligned}
R_{174} &={}-X_{109}X_{234} + X_{134}
\end{aligned}\]
\[\begin{aligned}
R_{175} &={}-X_{114}X_{235} + X_{134}
\end{aligned}\]
\[\begin{aligned}
R_{176} &={}-X_{127}X_{236} + X_{134}
\end{aligned}\]
\[\begin{aligned}
R_{177} &={}-2 - X_{105}^{2} - X_{153}^{2} - X_{162}^{2} + 2X_{169} - X_{169}^{2} - X_{175}^{2}\\
&\quad{} - X_{179}^{2} - X_{183}^{2} - X_{187}^{2} - X_{191}^{2} - X_{195}^{2} - X_{199}^{2} - X_{203}^{2}\\
&\quad{} - X_{207}^{2} - X_{211}^{2} + X_{215}^{2} - X_{237}^{2} - X_{238}^{2} - X_{239} - X_{239}^{2}
\end{aligned}\]
\[\begin{aligned}
R_{178} &={}-1 - 4X_{111}^{2} - 4X_{116}^{2} - 4X_{119}^{2} - 4X_{122}^{2} - 4X_{125}^{2} - 4X_{129}^{2}\\
&\quad{} - 4X_{132}^{2} + X_{134}^{2} - 4X_{136}^{2} - 4X_{140}^{2} - 4X_{143}^{2} - 4X_{146}^{2} - X_{240}^{2}\\
&\quad{} - X_{241}^{2} - X_{242} - X_{242}^{2} - 4X_{244}^{2}
\end{aligned}\]
\[\begin{aligned}
R_{179} &={}-X_{5}X_{243} + X_{134}
\end{aligned}\]
\[\begin{aligned}
R_{180} &={}X_{3}X_{134} - X_{6}X_{243} - X_{244}
\end{aligned}\]
\endgroup


\end{document}
